\section{Discussion}
\label{sec:discussion}

\paragraph{The gap result.}
Within each of the four judges, the same gap faults were caught less often when a record line was present, and
most of the lost catches became passes with no finding. Some mundane explanations can be set aside. The items are
identical across conditions. The drop also occurs when the line says only that no replacement is recorded, so it does
not depend on the record stating anything false. And the decoy test shows that most text-only catches were not
reflexive objections to older documents. What the design cannot set aside, the small and partly shared item sets and
the reused replies, is discussed in \S\ref{sec:limitations}.

We cannot say why the judges behaved this way, and at least three explanations fit what we observed. The judge may
simply be relying on the record as the task presents it. It is shown an ``amendment record lookup'', nothing suggests
the record is incomplete, and trusting it is defensible; on correct answers, following it removed most decoy false
alarms (\S\ref{sec:res-correct}). On that reading the failure belongs to a record that does not disclose its own gaps,
and what the judge lacks is any sign of registering the conflict with the texts. The line may instead be more salient
than the evidence, being short, explicit and placed just before the task, while the replacement has to be inferred
from two abstracts; we did not vary position or length. Or the judge may treat the status line as part of the question
rather than as evidence to weigh. A bare pass is compatible with all three, and it does not tell us whether the judge
noticed the conflict and resolved it in the record's favour or never noticed it.

\paragraph{Presentation.}
The presentation conditions were meant to separate these explanations, and they partly do. Rewording the gap did not
restore the catches for any judge, which lends no support to the simplest account, that judges are misled by an
explicit claim of currency, although 24 pairs cannot exclude a small effect of wording. For Gemini and GPT-5.5,
lowering the record's standing or warning that it was incomplete restored most or all of the lost catches, and Claude
recovered six of the eight it had lost, without reaching significance. This fits the first explanation: told that the
record might be wrong, these judges went back to the texts. It fits the salience explanation less well, since the note
and the warning leave the line as short, explicit and close to the task as before. For these three judges, however,
the decoy false alarms came back with the catches, to about their text-only level. The judges did not come to weigh
record against texts item by item; they moved from relying on one source to relying on the other. Gemma moved least,
catching at most five of the 24 gap faults however the damaged record was presented. Which judges respond in this way,
and why, is a question for more judges and more wordings than our one variant per condition.

Two related findings point in different directions. \citet{kapetanovic2026anchoring} found that a warning did not
reduce the total effect of judges' anchoring on prior scores, though it improved accuracy in one of their experiments.
Our warning differs in stating a specific and true defect of the record, with the contradicting evidence in view, and
either difference could matter. \citet{molfetta2026sycophants} report that larger models, reasoning models in particular, defer
more to authoritative but false legal citations. Our one open-weight judge responded least to the
presentation of the record, but the hosted judges' sizes are not published, and four judges could not test an effect
of size in any case.

\paragraph{Resemblance between documents.}
The decoy results contain two observations that should be kept apart. The judges did not flag older documents merely
for having a newer neighbour, since they caught far more replacements than they flagged updates. Yet their mistakes in
both directions tracked how alike the two documents looked. Given both texts, updates were flagged only when the titles were
near-identical, and Gemma rarely caught a replacement whose title had changed. Some of this reliance is reasonable,
since RFC replacements often do repeat their predecessor's title and abstract, and some flagged decoys are new
protocol versions that a reader might prefer to cite. The practical concern is transfer. An amended regulation may keep
its title, while a contract amendment can look nothing like the agreement it modifies. A judge evaluated on one kind
of collection may behave differently on the other, and nothing in a single-corpus evaluation, ours included, would show
it.

\paragraph{Scoring.}
Reason codes matter most on correct answers: under any-finding scoring, Claude, Gemini and Gemma would have appeared
to object to roughly two in five controls, nearly all for reasons unrelated to currency. They also make the gap result
measurable. A judge that calls a stale citation a ``content error'' has at least objected, and only reason-coded
scoring separates it from a judge that says nothing.
